\titledquestion{Test Paper Arrangement Task}

Midterm week is finally over! Now the CS101 course staff is filing the test papers. Suppose there are $t$ test papers waited to be filed, and the course staff was assigned $n$ file cabinets with various capacities. Mark the capacity of the file cabinet as $c_i$ ($i = 1, 2, \cdots, n$). It's guaranteed that $\sum_{i = 1}^n c_i \geq t$.\\



\begin{parts}

\part[5] Given the number of test papers $t$ and the capacities of each of the $n$ cabinets $c_1, c_2, \cdots, c_n$, come up with a greedy algorithm that can file the papers in cabinets such that the least number of cabinets is used. You should give out the steps of your algorithm \textbf{(as efficient as possible)} and the time complexity (tight).
\begin{solution}
    \begin{algorithm}[H]
        \caption{Greedy Algorithm for Test Paper Arrangement}
        \begin{algorithmic}[1]
            \State Sort the capacities of the cabinets in descending order: $c_1>c_2>\cdots>c_n$.
            \State Initialize the number of test papers left to be filed: $t$.
            \State Initialize the number of cabinets used: $m=0$.
            \For{$i=1$ to $n$}
                \State $t=t-c_i$
                \If{$t <= 0$}
                    \State $m=i$
                    \State \textbf{break}
                \EndIf
            \EndFor
            \State \Return $m$
        \end{algorithmic}
    \end{algorithm}
    \textbf{Time Complexity:} $O(n)$
\end{solution}

\newpage

\part[5] Prove the optimality of your algorithm using \textbf{Exchange Arguments} (any optimal solution can be iteratively exchanged to your greedy solution without making it worse). Answer the questions in step 4.5. to give your sufficient proof.\\
\textbf{Proof:}
\begin{enumerate}
    \item Assume $c_1>c_2>\cdots>c_n$ after sorted. Denote a solution $X=(x_1,x_2,\cdots,x_n)$ where we put $x_i$ test papers in the cabinet of capacity $c_i$. And the measurement $m(X)$ is counting how many cabinets have at least one test paper.
    \item Assume the optimal solution is $X^*$, and the greedy solution is $X$.
    \item Consider the first point of discrepancy between $X^*$ and $X$ at index $i$ (i.e. $x^*_i\ne x_i$).
    \item How to perform changes to make $x^*_i=x_i$? You should ensure that such changes is always valid in all possible scenarios of $X^*$.
    \begin{solution}
        If $x^*_i$ is greater than zero, which means it is used in $X^*$, we can randomly choose $x_i-x^*_i$ test papers from the rest of $X^*$, which is the set of $x^*_{i+1}$ to $x^*_n$. Then we put these test papers into the cabinet of capacity $c_i$. In that way, we can make $x^*_i=x_i$. This operation is always valid because sum of $X$ and $X^*$ are equal, so there are at least $x_i - x^*_i$ test papers in the rest of $X^*$.\\
        And for $x^*_i=0$, we can go on to find a j such that $x^*_j>0$ and $j>i$, then we can move $x^*_j$ test papers to the cabinet of capacity $c_i$. This operation is also valid because we sort the cabinets in descending order, so for $j>i$, $x^*_j \leq c_j<c_i$, which means the cabinet of capacity $c_i$ can hold the test papers from cabinet $j$.
    \end{solution}
    \item Why such changes doesn't make $X^*$ worse? You should consider all possible scenarios of $X^*$.
    \begin{solution}
        For situation where $x^*_i$ is greater than zero, the number of cabinets used in $X^*$ will not increase. Because we only move test papers from the rest of $X^*$ to a used cabinet of capacity $c_i$. That will do no contribution to the number of cabinets used.\\
        For situation where $x^*_i=0$, the number of cabinets used in $X^*$ will not increase. Even though we do let the cabinet of capacity $c_i$ hold some test papers, we also move all test papers from the cabinet of capacity $c_j$. So in general we are just exchanging the cabinets used in $X^*$, which will not increase the number of cabinets used.\\
        Thus we are sure to say that such changes doesn't make $X^*$ worse.
    \end{solution}
    \item Continue Applying steps 3.4.5. until $X^*=X$. Finally we prove that $X$ is an equally optimal solution.
\end{enumerate}

\end{parts}
\newpage